% !TEX root = ../Thesis.tex
% ============================================================
\chapter{Introduction} \label{chapter:introduction}
% ============================================================

% EDITORIAL NOTE (revision, 19 Aug 2026): the previous draft was flagged as
% repetitive. The cause was that four ideas each had three or four homes: the
% context-window argument, the top-k/expansion/path comparison, the description
% of SnapQuery, and the two-phase evaluation plan. This revision gives every
% idea exactly one home and cross-references it elsewhere:
%   * context-window argument ......... Section 1.2, once
%   * comparison of selection families  Chapter 2 only (was also in 1.2)
%   * what SnapQuery is ............... Section 2.7.4 only (was in 1.1, 1.3, 1.4, 2.7.4)
%   * two-phase evaluation ............ Section 1.4 only (was in Abstract, 1.4, 1.5)
%   * answer-node recall definition ... Section 1.3, once; used by name thereafter
% Section 1.2 now also states the measured properties of the graph, which
% replaces generic claims about clinical graphs with facts about this one.

A hospital knowledge graph and a large language model (LLM) are complementary in
an obvious way. The graph holds clinical facts in a form that can be inspected,
corrected and audited one edge at a time; the model lets someone ask about those
facts without knowing a query language. Putting the two together is attractive
precisely because neither substitutes for the other: the graph supplies
trustworthiness, the model supplies access.

What decides whether such a system works is a step that is easy to underestimate.
Before an LLM can answer anything, something must choose which part of the graph
it gets to see. This thesis studies that choice, called \emph{subgraph
selection}, on a clinical knowledge graph that is in daily use at University
Hospital Basel, and compares one principled strategy for making it against the
system that answers these questions there today.

\section{Motivation} \label{sec:motivation}

Two things make this worth studying rather than assuming.

The first is that the leading method in the general-purpose setting has never
been tested in this one. G-Retriever~\cite{he2024gretriever} formulates subgraph
selection as a Prize-Collecting Steiner Tree (PCST) problem, so that relevance
and subgraph size are traded off inside a single optimisation rather than by
post-hoc filtering, and reports large gains on general textual graphs. Its
benchmarks, however, are small and open-domain: on WebQSP each question arrives
with its own pre-extracted subgraph averaging 1{,}371 nodes. A hospital graph is
a different object, and Section~\ref{sec:problem-statement} gives the measured
respects in which it differs.

The second is the nature of the comparison. The system this thesis measures
against, SnapQuery (Section~\ref{subsec:snapquery-background}), is not a research
baseline reimplemented for the purpose. It is a deployed service that clinicians
and data managers at the STCS Data Center rely on, built by the people who
maintain the graph, and tuned against the questions its users actually ask.
Outperforming it is a stronger claim than outperforming a paper baseline, and
failing to outperform it is a more informative result.

\section{Problem Statement} \label{sec:problem-statement}

Passing the whole graph to the model is not an option. Serialising a graph into
text is straightforward in principle, but LLM context windows are finite and a
hospital-scale graph produces far more text than any current model can read at
once, at a cost in time and compute that would be impractical regardless. So
something has to select, and selection has two failure modes: retrieve too little
and the facts an answer depends on are missing; retrieve too much and the result
is expensive to process, hard to read, and diluted by material that competes with
the relevant part for the model's attention. Chapter~\ref{chapter:background}
compares the families of strategies that navigate this trade-off; what matters
here is that the trade-off exists and that no strategy escapes it.

Three properties of the specific graph studied here sharpen the problem
considerably. All three were measured during this project, and none of them holds
in the benchmarks G-Retriever was evaluated on.

% TODO: confirm with the STCS Data Center that reporting these aggregate counts
% in a public thesis is acceptable. They are cohort-level statistics with no
% patient-level information, but the confirmation should be on record.
\paragraph{Scale.} The graph holds 21{,}537{,}305 nodes and 72{,}218{,}890 edges
for 1{,}197 patients, roughly 18{,}000 nodes per patient. Selecting over the
whole graph per question is therefore not merely expensive but structurally
different from the published setting, where the candidate region was fixed in
advance and small.

\paragraph{Sparse and shared text.} Almost no node carries text of its own.
Laboratory results, which with their associated samples and test records make up
roughly 95\% of the graph, are described only through the terminology codes they
reference, and 4{,}832{,}744 laboratory tests resolve to just 544 distinct LOINC
descriptions. A similarity score computed over that text can take only 544
distinct values across millions of nodes, so it can identify which
\emph{analyte} a question is about but not which \emph{measurement}.

\paragraph{Degenerate connectivity.} A handful of administrative nodes dominate
the graph's structure: 168 provenance nodes have an average degree of 69{,}525,
and six anatomical-site nodes absorb 3{,}454{,}881 edges. Any two nodes are two
or three hops apart through such a node, which makes an unqualified
connectivity requirement nearly vacuous.

Finally, and independently of the graph, the relevance signal itself is only a
proxy. Cosine similarity between a question and a node's text is not relevance,
and terse clinical labels are not the open-domain prose most sentence encoders
were trained on. A node can rank highly for the wrong reason, or poorly despite
being exactly right, because it uses a clinically equivalent but differently
worded term. A selection strategy has to remain useful under an imperfect signal,
not only under a favourable graph.

\section{Research Questions} \label{sec:research-objectives}

The thesis addresses three questions.

\begin{description}
\item[RQ1] Does PCST-based subgraph selection retrieve answer entities more
effectively than SnapQuery on the same clinical questions, measured as
\emph{answer-node recall at a given subgraph size}: the fraction of entities
constituting a correct answer that appear in the retrieved subgraph, reported
against the size of that subgraph rather than as a single number?

\item[RQ2] How much of any difference is attributable to the connectivity
constraint itself, rather than to the underlying similarity scores? This is
answered by internal ablations that hold the similarity signal fixed and vary
only how connectivity is handled: similarity-only top-$k$ selection, and
breadth-first expansion from seed nodes.

\item[RQ3] Does augmenting the PCST prize function with \emph{intent} signals
(recognising that a question concerns a treatment, a laboratory value, or a
complication) and \emph{schema} signals (raising the prizes of nodes whose SPHN
type matches that intent) improve the trade-off further?
\end{description}

The working hypothesis for RQ1 is that PCST-based selection achieves higher
answer-node recall at a given subgraph size, and that the margin is largest on
questions requiring several entities to be connected through the graph rather
than a single node to be matched. The reasoning is asymmetric in an interesting
way: SnapQuery returns exactly the rows its generated query names, so an entity
it does not explicitly ask for cannot appear in its result, whereas the
connectivity term in PCST admits bridge nodes and thereby gives indirectly
related answers a chance to surface. RQ3 is pursued only if the work on RQ1 and
RQ2 stays on schedule.

\section{Scope and Delimitations} \label{sec:thesis-scope}

The evaluation proceeds in two phases. Phase~1, the core of the thesis, measures
retrieval quality alone: whether the retrieved subgraph contains the answer
entities, and how large it is, independently of whether a downstream model would
phrase a fluent answer from it. Phase~2 measures answer quality: an LLM generates
a natural-language answer from each retrieved subgraph, which is then assessed
for faithfulness and hallucination~\cite{huang2023hallucination}. Phase~2 belongs
to this thesis rather than to future work.

Query latency is deliberately excluded. It is a property of an implementation
and a deployment rather than of a selection strategy, and including it would
reward engineering effort instead of the retrieval behaviour under study.

The graph is the Neo4j graph that SnapQuery itself queries, scoped to the STCS
transplant cohort and structured according to the SPHN
schema~\cite{sphnschema}. The schema is public and versioned; the instance data
is governed clinical data that remains on University Hospital Basel and BioMedIT
infrastructure throughout. Development therefore proceeds against the schema,
against synthetic graphs reproducing its measured structure, and against the real
data once access permits, with final validation on the real cohort.

\section{Contributions} \label{sec:contributions}

\begin{itemize}
\item An implementation of PCST-based subgraph selection for an SPHN-structured
clinical knowledge graph, verified to reproduce the published G-Retriever
implementation exactly on shared inputs, so that any difference in behaviour on
clinical data is attributable to the data rather than to the reimplementation.

\item A \emph{graph preparation} procedure that makes an SPHN instance graph
amenable to text-based retrieval at all: folding shared terminology nodes into
the text of the nodes that reference them, pruning provenance hubs, and scoping
selection to a patient or cohort. Each step is stated as a methodological choice
with a measured justification, and each can be disabled for ablation.

\item Three findings about how G-Retriever transfers to this setting, each
established both analytically and empirically:
  \begin{enumerate}
  \item its edge-prize mechanism collapses when a graph has few, heavily shared
  relation types, because the per-tier prize budget is divided among all
  identically-scored edges; under the published default parameters this drives
  the effective per-edge cost to a negligible value and the method returns
  essentially the entire candidate region;
  \item with text shared at high multiplicity, the top-$k$ node selection is
  determined by tie-breaking rather than by similarity, so two implementations
  of the same algorithm return different subgraphs for the same question;
  \item resolving terminology codes to canonical descriptions, rather than
  embedding the labels the graph happens to carry, is what makes the retrieval
  signal usable, and it is also what makes results portable across the
  multilingual terminologies an SPHN graph combines.
  \end{enumerate}

\item An evaluation protocol that compares a subgraph retriever against a
production query-generating agent on a shared metric, despite the two exposing
different artefacts, together with an explicit treatment of the questions this
metric cannot score, so that they are reported rather than silently counted as
failures.

\item A human-validated gold set of approximately 50 questions with manually
confirmed answer entities, used to establish correctness independently of
SnapQuery's own output.

\item An assessment of answer faithfulness and hallucination for answers
generated from the retrieved subgraphs (Phase~2), and, as an optional extension,
of the intent- and schema-aware prize function of RQ3.
\end{itemize}

\section{Thesis Structure} \label{sec:thesis-structure}

Chapter~\ref{chapter:background} introduces the technical background and
positions PCST-based selection among the alternatives. Chapter~3 gives the
mathematical treatment of G-Retriever and PCST that
Chapter~\ref{chapter:background} deliberately leaves out. Chapter~4 sets out the
methodology: the graph preparation procedure, the strategies compared, the
evaluation protocol and the question sets. Chapter~5 describes the
implementation. Chapter~6 reports and discusses the Phase~1 and Phase~2 results,
and Chapter~7 concludes.
